\section{Harmonic models and row characters}
\label{sec:preliminaries}

For the remainder of the paper, except where a more general base is
specified, $K$ is a field of characteristic zero.  We retain the
polynomial row convention and the notation $P_{n,k}$ and $R_{n,k}$
from the introduction.

\subsection{Harmonics and the operator module}

For $\alpha=(\alpha_1,\ldots,\alpha_k)\in\N^k$, define
\[
 D_\alpha=\sum_{i=1}^n
       \prod_{a=1}^k\partial_{x_i^{(a)}}^{\alpha_a}.
\]
The full diagonal harmonic space is
\begin{equation}\label{eq:harmonic-definition}
 \Harm_{n,k}=\bigcap_{\alpha\ne0}\ker D_\alpha.
\end{equation}
Polarized power sums generate $P_{n,k}^{S_n}$ in characteristic
zero.  Their apolar equations identify
\eqref{eq:harmonic-definition} with the orthogonal complement of
the Hilbert ideal for the Fischer pairing
\[
 \langle f,g\rangle=f(\partial)g\big|_{X=0}.
\]
Over the reals this pairing is positive definite on each
homogeneous component.  The quotient projection therefore
restricts to an isomorphism
\begin{equation}\label{eq:harmonic-quotient}
 \Harm_{n,k}\xrightarrow{\ \sim\ }R_{n,k}.
\end{equation}
The spaces and maps are defined over $\Q$, so the assertion
holds over $K$ by extension of scalars.  The differential
equations are stable under row changes; consequently the
isomorphism is $\GL_k\times S_n$-equivariant with our
polynomial row convention.

The ordinary Vandermonde is harmonic.  Partial derivatives
preserve \eqref{eq:harmonic-definition}, and
\begin{equation}\label{eq:polarization-commutator}
 [D_\alpha,E_{a,b}^{(s)}]
       =\alpha_a D_{\alpha-e_a+s e_b},
\end{equation}
where a zero coefficient gives zero.  This proves
$\Pol_{n,k}\subseteq\Harm_{n,k}$.
The operators with $s=1$ give the row Lie algebra.  As every
homogeneous part is a polynomial representation, $\Pol_{n,k}$
is a $\GL_k\times S_n$-submodule.  It is graded: the initial
Vandermonde is homogeneous and all generating operators are
homogeneous for total degree.

The linear harmonic equations say that every harmonic is
translation-invariant in each row.  Write
\[
 W_n=\{(z_1,\ldots,z_n)\in K^n:\ \textstyle\sum_i z_i=0\}.
\]
Using the invariant dot product to identify $W_n$ with its
dual, we can regard $\Harm_{n,k}$ and $\Pol_{n,k}$ as
derivative-closed submodules of $\Sym(K^k\otimes W_n)$.
In particular every occurring row partition has length at
most $n-1$ by the Cauchy decomposition.

\subsection{Stable row coefficients}

\begin{lemma}[Trailing-zero stability]\label{lem:stability}
For $\ell(\lambda)\leq k$, the highest-weight multiplicity
space of row weight $\lambda$ in $\Harm_{n,k}$ is naturally
the highest-weight multiplicity space of
$(\lambda,0)$ in $\Harm_{n,k+1}$.
\end{lemma}

\begin{proof}
A polynomial with last row degree zero does not use the last
row.  Every differential equation involving that row vanishes
on it.  The other equations are exactly the equations in
$k$ rows.  The additional row raising operators also differentiate
in the last row and hence vanish.  The remaining highest-weight
conditions are unchanged.
\end{proof}

Thus adjoining an unused row preserves every full harmonic
multiplicity in \eqref{eq:character-decomposition}.  The stable
character is finite: its row lengths
are bounded by $n-1$, and $R_{n,n-1}$ is finite-dimensional:
the invariant ideal contains the ordinary invariant ideal in
each row, so the quotient is a quotient of a tensor product
of ordinary coinvariant rings.
We write its label coefficients as symmetric functions
\[
 C_\mu^{(n)}=C_\mu^\Harm
       =\sum_\lambda c_{\lambda,\mu}^{(n)}s_\lambda,
 \qquad A_n=A_n^\Harm=C_{(1^n)}^{(n)}.
\]
Evaluation in $k$ variables recovers the character in $k$ row
alphabets.  For symmetric functions, $f^\perp$ denotes the
adjoint of multiplication by $f$ for the Hall inner product, in
which the Schur functions are orthonormal.

For $\Pol$, we use the direct system obtained by adjoining unused
rows.  For a fixed row partition its highest-weight spaces form
an increasing sequence inside the fixed finite-dimensional
full-harmonic space of Lemma~\ref{lem:stability}; the sequence
therefore stabilizes.  This defines the stable characters
\[
 C_\mu^{\Pol}=\sum_\lambda c_{\lambda,\mu}^{\Pol}s_\lambda,
 \qquad A_n^{\Pol}=C_{(1^n)}^{\Pol}.
\]
The commuting normal form in
Theorem~\ref{thm:r3-normal-form} will identify these stable
characters with their actual specialization in every finite
number of rows.

\subsection{A parabolic invariant quotient}

For a set of labels $T\subseteq[n]$, write $S_T$ for the
permutations supported on $T$.  Set
$X_i=(x_i^{(1)},\ldots,x_i^{(k)})$.
Passing to tail invariants leaves only the selected labels:
every tail power sum becomes the negative of the selected power
sum.  This elementary elimination controls row length throughout
the paper.

\begin{lemma}\label{lem:parabolic}
For $0\leq r\leq n$, the natural map
\[
 K[X_1,\ldots,X_r]\longrightarrow R_{n,k}^{S_{\{r+1,\ldots,n\}}}
\]
is surjective.  Consequently every row partition occurring
in the target has length at most $r$.
\end{lemma}

\begin{proof}
For $r\ge1$, surjectivity is the polynomial case of
Theorem~\ref{thm:universal-selected-trace-presentation}; for $r=0$
it is $R_{n,k}^{S_n}=K$.  Finally,
\[
 K[X_1,\ldots,X_r]
  =\Sym(K^k\otimes K^r)
  \cong\bigoplus_\lambda
       \Schur_\lambda K^k\otimes\Schur_\lambda K^r,
\]
and $\Schur_\lambda K^r=0$ for $\ell(\lambda)>r$.
\end{proof}

\begin{proposition}[The length upper bound]\label{prop:length-upper}
If $c_{\lambda,\mu}^{(n)}\ne0$, then
\[
 \ell(\lambda)\leq n-\mu_1.
\]
The same bound holds in every $\Pol_{n,k}$.
\end{proposition}

\begin{proof}
The restriction of $S^\mu$ to a subgroup $S_{\mu_1}$ contains
the trivial representation.  By the branching rule, one may
remove every box below the first row, in a sequence of removable
corners, and reach the one-row partition $(\mu_1)$.
Thus an occurrence of $\Schur_\lambda K^k\otimes S^\mu$
gives an occurrence of $\Schur_\lambda K^k$ in
$R_{n,k}^{S_{\mu_1}}$.
Apply Lemma~\ref{lem:parabolic} with $r=n-\mu_1$.
The assertion for $\Pol$ follows from its inclusion in
$\Harm\cong R$.
\end{proof}

